% WORKING MATHEMATICAL DRAFT: sharp factor, real minimizer, and complex
% minimizer are not Lean-verified. The projection lower bound and exact
% four-dimensional trace ratio are already Lean-verified. Keep the mixed-status
% section behind \includeUnformalizedDisplacementtrue until the other proofs compile.
\section{Quantitative bounds for full displacement (working draft)}
\label{sec:unformalized-sharp-displacement}

\begin{tcolorbox}[
  enhanced, breakable, colback=white, colframe=red!80!black,
  colbacktitle=red!6!white, coltitle=red!80!black,
  title={WORKING EXTENSION --- sharp optimality theorems not yet Lean-verified},
  fonttitle=\bfseries, boxrule=1.05pt,
  left=9pt, right=9pt, top=9pt, bottom=9pt,
  before skip=12pt, after skip=12pt]

The projection-distance lower bound below is verified in Lean for every
unitarily invariant seminorm over real or complex finite-dimensional spaces.
The exact trace-ratio identity for the parameterized real counterexample is
also verified. The universal sharp multiplicative bound, global trace
optimality of the real competitor, and the complex-space trace optimum have
mathematical proofs here but remain outside the compiled Lean development.
The verified Proposition~4.4 refutation is
\cref{thm:small-angle,cor:no-threshold}.

Let $\mathbb F$ denote $\mathbb R$ or $\mathbb C$, and let $U,V$ be subspaces
of equal dimension $k$ in an $n$-dimensional $\mathbb F$-Hilbert space.
Write $P=P_U$ and $Q=P_V$ for their orthogonal projections and
$0\leq\theta_1\leq\cdots\leq\theta_k\leq\pi/2$ for their
principal angles, including any zero angles. A norm $N$ on square operators
is \emph{unitarily invariant} when $N(AXB)=N(X)$ for orthogonal or unitary
$A,B$. The notation $\|\cdot\|_1$ denotes the trace norm.

\begin{proposition}[Projection-distance lower bound]
\label{prop:draft-projection-lower}
For every orthogonal or unitary map $W$ satisfying $WU=V$, and every
unitarily invariant norm $N$,
\begin{equation}
  N(P-Q)\leq N(I-W).
  \label{eq:draft-projection-lower}
\end{equation}
\end{proposition}

\begin{proof}
Decompose $W$ with respect to $U\oplus U^\perp$:
\[
 W=\begin{pmatrix}A&B\\ C&D\end{pmatrix},
 \qquad J=2P-I=\begin{pmatrix}I&0\\0&-I\end{pmatrix}.
\]
For $K=I-W$, its off-diagonal part is
\begin{equation}
 \frac{K-JKJ}{2}=-\begin{pmatrix}0&B\\C&0\end{pmatrix}.
 \label{eq:draft-offdiag}
\end{equation}
By the triangle inequality and invariance under conjugation by $J$,
\[
 N\!\left(\frac{K-JKJ}{2}\right)
 \leq\frac{N(K)+N(JKJ)}2=N(K).
\]
The singular values of $A=P W|_U$ are $\cos\theta_j$: indeed,
$W|_U$ is an isometry onto $V$, so these are the singular values of
$P|_V$ in an orthonormal frame. The identities from $W^*W=WW^*=I$ give
\[
 C^*C=I_U-A^*A,
 \qquad BB^*=I_U-AA^*.
\]
Thus the nonzero singular values of each of $B$ and $C$ are the
$\sin\theta_j$, with the appropriate zeros. The off-diagonal operator in
\cref{eq:draft-offdiag} has singular values $\sin\theta_j$, each twice,
together with zeros. The cosine-sine decomposition for a pair of projections
shows that $P-Q$ has exactly the same singular values
\citep{stewart1990matrix,bhatia1997matrix}.
Unitary invariance therefore identifies the left side of the preceding
pinching estimate with $N(P-Q)$, proving \cref{eq:draft-projection-lower}.
\end{proof}

This lower bound is independently verified in Lean as
\texttt{projectionDifference\_uiNorm\_le\_displacement}, with the stronger
seminorm hypothesis and without an acute-angle restriction.

\begin{theorem}[Sharp multiplicative displacement estimate]
\label{thm:draft-sharp-displacement}
Assume that $\alpha=\theta_k<\pi/2$, and let $R$ be the canonical direct
rotation from $U$ to $V$. For every unitarily invariant norm $N$,
\begin{equation}
 N(I-R)\leq\sec(\alpha/2)
 \min_{W:\,WU=V}N(I-W).
 \label{eq:draft-sharp-displacement}
\end{equation}
For each fixed $0<\alpha<\pi/2$, the factor $\sec(\alpha/2)$ is
attained by a real four-dimensional pair and the trace norm. In particular,
under the angle hypothesis $\alpha\leq\pi/3$ the sharp common factor is
$2/\sqrt3$; uniformly for acute pairs, $\sqrt2$ is a valid upper bound.
\end{theorem}

\begin{proof}
In the principal-plane decomposition, the canonical direct rotation acts
on the plane of angle $\theta_j$ as a rotation through that angle and is
the identity on the remaining orthogonal complement. The nonzero singular
values of $I-R$ are consequently
\[
 2\sin(\theta_j/2),\quad 2\sin(\theta_j/2)
 \qquad (j=1,\ldots,k),
\]
with zeros for the remaining dimensions. Similarly, the singular values of
$P-Q$ are $\sin\theta_j$, each twice, with zeros. For every $j$,
\[
 2\sin(\theta_j/2)
 =\frac{\sin\theta_j}{\cos(\theta_j/2)}
 \leq \sec(\alpha/2)\sin\theta_j.
\]
Symmetric-gauge monotonicity of unitarily invariant norms
\citep{bhatia1997matrix} yields
\[
 N(I-R)\leq\sec(\alpha/2)N(P-Q).
\]
Combining this with \cref{eq:draft-projection-lower}, valid for every
admissible $W$, proves \cref{eq:draft-sharp-displacement}. The minimum
exists because the admissible orthogonal or unitary matrices form a closed
subset of a compact group.

For the real family in \cref{thm:small-angle}, the two principal angles are
$\theta_1=\theta_2=\alpha$. Hence $\|P-Q\|_1=4\sin\alpha$, while the
verified formulas give
\[
 \|I-W_\alpha\|_1=4\sin\alpha,
 \qquad \|I-R_\alpha\|_1=8\sin(\alpha/2).
\]
The projection lower bound proves that $W_\alpha$ is a \emph{global}
trace-norm minimizer for this pair. The ratio of the canonical displacement
to that minimum equals $\sec(\alpha/2)$, giving equality in
\cref{eq:draft-sharp-displacement} for every fixed $0<\alpha<\pi/2$.
The exact identity $\cos(\alpha/2)\|I-R_\alpha\|_1
=\|I-W_\alpha\|_1$ is separately Lean-verified as
\texttt{SmallAngleCounterexample.cos\_half\_mul\_directTrace\_eq\_competitorTrace};
the assertion that $W_\alpha$ attains the global minimum remains in the
working mathematical draft.
At $\alpha=\pi/3$ this ratio is $2/\sqrt3$, and as
$\alpha\uparrow\pi/2$ it approaches $\sqrt2$.
\end{proof}

\begin{proposition}[Exact trace-norm optimum over complex spaces]
\label{prop:draft-complex-trace-optimum}
If $U,V$ are equal-dimensional subspaces of a finite-dimensional complex
Hilbert space, then
\begin{equation}
  \min_{W:\,WU=V}\|I-W\|_1
  =\|P-Q\|_1
  =2\sum_{j=1}^k\sin\theta_j.
  \label{eq:draft-complex-trace-optimum}
\end{equation}
No acuteness hypothesis is needed.
\end{proposition}

\begin{proof}
The lower bound is \cref{prop:draft-projection-lower}. The principal-angle
decomposition gives mutually orthogonal two-dimensional complex planes
for the nonzero angles and a common fixed subspace for the zero angles
and unused ambient dimensions. On one such plane use the unitary rotations
\begin{equation}
 R_\theta=\begin{pmatrix}
   \cos\theta&i\sin\theta\\ i\sin\theta&\cos\theta
 \end{pmatrix},
 \qquad W_\theta=e^{i\theta}R_\theta.
 \label{eq:draft-complex-phase}
\end{equation}
For the source line $\operatorname{span}_{\mathbb C}\{(1,0)\}$,
these maps have the same target line, since multiplication by a scalar phase
does not alter a complex subspace. The eigenvalues of $R_\theta$ are
$e^{i\theta},e^{-i\theta}$, while those of $W_\theta$ are
$1,e^{2i\theta}$. Since both are normal,
\[
  \|I-R_\theta\|_1=4\sin(\theta/2),
  \qquad \|I-W_\theta\|_1=2\sin\theta.
\]
This construction is valid also for $\theta=\pi/2$, with the two-dimensional
plane obtained by choosing orthonormal vectors in the orthogonal subspaces.
Take the orthogonal direct sum of the maps $W_{\theta_j}$ over all nonzero
principal angles and the identity on the remaining summand. The resulting
unitary carries $U$ onto $V$ and has trace-norm displacement
$2\sum_j\sin\theta_j=\|P-Q\|_1$, attaining the lower bound.
\end{proof}

\paragraph{The real four-dimensional realization.}
Identifying $\mathbb C^2$ with $\mathbb R^4$ and complex multiplication by
$i$ with $J=\left(\begin{smallmatrix}0&-1\\1&0\end{smallmatrix}\right)$
turns \cref{eq:draft-complex-phase} into the real block matrices in
\cref{eq:wtheta}, including the canonical rotation $R_\theta$. Complex
singular values are duplicated under realification, producing the trace
norms in \cref{eq:competitor-trace,eq:direct-trace}. Davis and Kahan already
used a complex scalar phase to exhibit trace-norm nonminimality
\citep[Example~4.2]{davis1970rotation}; the realification explains why
their subsequent restriction to real spaces does not restore the asserted
minimality. Bolotnikov, Li, and Rodman later studied this norm-minimization
problem and the role of $Q$-norms in detail
\citep{bolotnikov2001minimal}.

\end{tcolorbox}
